\documentclass[11pt]{article}
\usepackage[margin=1in]{geometry}
\usepackage{amsmath,amssymb,booktabs,authblk,url}
\newcommand{\Ca}{C\ensuremath{\alpha}}
\newcommand{\Cfp}{C4\ensuremath{'}}

\title{\textbf{Which Carbon Takes the Proton?}\\[0.3em]
\large A parameter-free computation of the branch that costs GAD65 its cofactor}
\author[1]{Claes Johnson}
\affil[1]{KTH Royal Institute of Technology, Stockholm, Sweden \\ \texttt{claesjohnson@gmail.com}}
\date{\today \\[0.4em]\small\textit{Developed in collaboration with Claude (Anthropic).}}

\begin{document}
\maketitle

\begin{abstract}
\noindent Glutamate decarboxylase GAD65 --- the major type-1-diabetes autoantigen --- destroys its
own cofactor as a side reaction of its normal catalytic cycle. After CO$_2$ departs, the resulting
quinonoid intermediate can be reprotonated at either of two carbons two bond lengths apart: at
\Ca{}, which releases GABA and hands back an intact pyridoxal-5$'$-phosphate, or at \Cfp{}, which
gives pyridoxamine phosphate, falls out of the active site, and leaves the enzyme in the
cofactor-free form implicated in its autoantigenicity. Which carbon is favoured therefore sets the
rate of self-inactivation.

We compute that preference with RealQM --- a real-space method with no basis set, no
exchange--correlation functional and nothing fitted --- on an active site taken from the GAD65
crystal structure, with the substrate grafted where the catalytic lysine vacates and the
stereoelectronic alignment read off the measured ring normal rather than imposed. A bare proton is
placed at each site in turn. The two configurations carry the same twenty-seven kernels, the same
sixty-three electrons and the same grid, and differ in the position of one atom, so that the large
common errors cancel; this is the only construction in which the present framework's energies are
trustworthy, and we report nothing that falls outside it.

The preference is for \Cfp{} --- the destructive branch --- and it is electronic in origin: the
proton there pays $1.15$~Ha more in kernel repulsion and recovers $2.29$~Ha more in electronic
attraction. The direction is stable under refinement; the magnitude is not, and no magnitude is
claimed. The result is consistent with solution chemistry that the calculation was not tuned to
reproduce: free PLP with an amino acid transaminates rather than decarboxylates. On this evidence
the bare quinonoid prefers to destroy its cofactor, and catalysis is what an enzyme must impose
against that preference.
\end{abstract}

\section{The question}

Glutamate decarboxylase converts glutamate to GABA using pyridoxal-5$'$-phosphate. In the accepted
mechanism the substrate forms an external aldimine with the cofactor, the \Ca--COO$^-$ bond breaks
in the orientation perpendicular to the ring that Dunathan's rule selects, and the electron pair
left behind delocalises into the protonated pyridine ring: the quinonoid intermediate, an object
bright enough at $500$~nm to be followed in stopped-flow.

What happens next is a fork, and it is the subject of this article. A proton arrives, and the
conjugated chain offers it two carbons:
\begin{center}
\small
\begin{tabular}{@{}lll@{}}
\toprule
protonated at & gives & outcome\\
\midrule
\Ca{} --- the substrate's carbon & product aldimine & GABA released, PLP regenerated\\
\Cfp{} --- the cofactor's carbon & ketimine $\to$ PMP & cofactor lost, enzyme left \emph{apo}\\
\bottomrule
\end{tabular}
\end{center}
\noindent They are $2.37$~\AA{} apart on the same $\pi$ system. One returns the enzyme to service;
the other converts its cofactor into something that dissociates.

This matters beyond enzymology because GAD65 is the major type-1-diabetes autoantigen, and its
propensity to lose PLP and populate a flexible apo form has been linked to that autoantigenicity.
The chain runs: branch preference $\to$ rate of cofactor loss $\to$ size of the apo pool $\to$ (on
the standing hypothesis) autoantigenicity. We address the first link only. The later ones belong to
enzymology and immunology, and the hypothesis joining them is not ours to test.

\section{What can be computed, and what cannot}

RealQM represents an $N$-electron system by $N$ real-valued wave functions on non-overlapping
supports in three-dimensional space, with the physics in one functional --- Coulomb potential energy
without self-repulsion, plus kinetic energy --- and nothing fitted. It has no basis set and no
exchange--correlation functional, so there is no choice to make that could steer the answer, and no
population-analysis convention is needed to say where charge sits: the supports \emph{are} the
partition.

Against that stands a limitation established by calibration and worth stating before any result.
This framework's \emph{absolute} energetics are unreliable. Its proton affinity for ammonia comes
out at $-58$~kcal/mol against $+204$ measured --- the wrong sign --- and stays wrong across nine
variants of how the nitrogen is represented. The reason is arithmetic rather than mysterious: a
proton affinity differences two systems of \emph{different composition}, one of which has an extra
kernel, so the errors do not cancel. The same objection retires the obvious design for the present
problem, in which one compares a cofactor-bearing system with a cofactor-free one: those are systems
of twenty-four and thirteen atoms, and no amount of refinement makes their difference meaningful.

What does survive is a difference \emph{within one composition}. Two configurations of the same
atoms, same electron count, same kernels and same grid share their errors, and the errors subtract.
The quantity this article reports is of that kind and only that kind.

\section{The calculation}

\paragraph{Geometry from the structure.} The active site is taken from the GAD65 crystal structure
\cite{fenalti}, which contains the cofactor and the catalytic lysine as a single covalent residue,
LLP396, together with the aspartate that hydrogen-bonds the ring nitrogen at $2.46$~\AA{} and a bound
molecule of the product. The substrate is grafted at the positions the lysine vacates on forming the
external aldimine --- imine nitrogen at the crystallographic NZ, \Ca{} at CE --- and the breaking
\Ca--COO bond is laid along the \emph{measured} ring normal, so that the Dunathan alignment comes out
of the structure rather than being imposed: it checks to $0.0^\circ$. Bond lengths in the grafted
fragment agree with standard values (C4$'$--N7 $1.34$, N7--\Ca{} $1.49$~\AA). Hydrogens are added,
X-ray seeing none. The nuclei are held throughout; only the charge relaxes.

\paragraph{The probe.} A bare nucleus --- no electrons of its own, a unit kernel --- is a proton, and
is how proton transfer has been represented in this framework before. It is placed $2.00$~\AA{} from
\Ca{} on the face that CO$_2$ vacates, and in the second configuration $2.00$~\AA{} from \Cfp{} on
the face opposite its own hydrogen. The geometry is symmetric between the two: in each the probe is
$2.00$~\AA{} from its host and $3.88$~\AA{} from the other carbon, so neither site is privileged by
its placement.

\paragraph{What is held identical.} Both configurations carry twenty-seven kernels, sixty-three
electrons, total kernel charge sixty-four, one box and one mesh. Position by position, exactly one
atom differs between them. The reported quantity is
\[
\Delta = E(\text{proton at } \Ca{}) - E(\text{proton at } \Cfp{}),
\]
and positive $\Delta$ means \Cfp{} is preferred.

\section{Result}

\begin{center}
\small
\begin{tabular}{@{}rrrrl@{}}
\toprule
steps & $E(\Ca)$ (Ha) & $E(\Cfp)$ (Ha) & $\Delta$ (Ha) & residuals\\
\midrule
$2000$ & $-6.968$ & $-8.208$ & $+1.241$ & $2.08$ / $2.04$\\
$4000$ & $-7.364$ & $-8.499$ & $+1.136$ & $0.52$ / $0.53$\\
\bottomrule
\end{tabular}
\end{center}

\noindent $\Delta$ is positive at both step counts: \textbf{the proton prefers \Cfp{}}, the branch
that destroys the cofactor.

\paragraph{The preference is electronic, and it wins against the nuclei.} The kernel--kernel energy
the proton pays is computable in closed form from the geometry, and it runs the other way: $6.98$~Ha
at \Ca{} against $8.13$~Ha at \Cfp{}, because \Cfp{} sits nearer the ring and so nearer more kernels.
Subtracting,
\begin{center}
\small
\begin{tabular}{@{}lr@{}}
\toprule
contribution & Ha\\
\midrule
kernel--kernel, favouring \Ca{} & $-1.15$\\
electronic, favouring \Cfp{} & $+2.29$\\
\midrule
net, favouring \Cfp{} & $+1.14$\\
\bottomrule
\end{tabular}
\end{center}
\noindent So the result is not an electrostatic accident of where the nuclei sit. The proton is drawn
to \Cfp{} by the electrons, against a nuclear term that opposes it by half as much.

\paragraph{How far the cancellation goes, measured rather than assumed.} Between the two step counts
the individual energies moved by $0.40$ and $0.29$~Ha while $\Delta$ moved by $0.105$ --- a
cancellation of $73.5$ per cent. That is good enough to fix the sign, the drift being $66$~kcal/mol
against a difference of $713$, and not good enough to fix the value. It is also markedly worse than
the $99.5$ per cent obtainable when two configurations are related by an exact symmetry, which
\Ca{} and \Cfp{} are not. \textbf{We therefore report a direction and a mechanism, and no number.}

\paragraph{A consistency check the calculation was not tuned to pass.} PLP and an amino acid in
solution, with no enzyme present, undergo transamination readily --- which is protonation at \Cfp{}.
If the intrinsic preference were for \Ca{}, free PLP would decarboxylate instead. The computed
preference agrees with the solution chemistry, and nothing in the model was adjusted to make it do
so.

\section{Scope}

\begin{itemize}
\item \textbf{A direction, not a rate.} A branch \emph{ratio} requires barriers and a rate theory. A
minimisation over static densities contains no time and can supply neither. Nothing here predicts how
often GAD65 inactivates itself.
\item \textbf{One elementary step.} The catalytic cycle has six. Transaldimination, reprotonation and
product release are not modelled.
\item \textbf{Almost no enzyme.} The crystal supplies the geometry, and Asp364 is present, but there
is no protein beyond it, no solvent, and --- importantly --- no Lys396, which is the actual proton
donor. The enzyme's control over this branch is exerted by \emph{positioning} that donor, and a free
proton measures the intrinsic preference the enzyme works against, not the enzymatic outcome.
\item \textbf{Not a first.} Quantum-mechanical/molecular-mechanical work has treated the equivalent
\Ca{}-to-\Cfp{} transfer in a related PLP enzyme, with explicit lysine shuttle and free-energy
profiles, which is more than is offered here. What is new is the absence of fitted content: no
functional was chosen and no population convention was needed.
\end{itemize}

\section{What follows}

The natural next step is the comparison this article does not yet make. GAD67, the paralogue, sheds
its cofactor far more rarely, and its crystal structure was solved alongside GAD65's. The two active
sites are chemically identical --- the same aspartate, histidine, leucine, threonine, asparagine and
alanine, with the numbering shifted by nine --- so a comparison between them is a difference of
\emph{geometry at fixed composition}, which is the construction this framework supports. The same
two-position measurement can therefore be repeated in the GAD67 site and the two preferences set side
by side.

Both outcomes would be informative, which is what makes it worth doing. A shifted preference would
locate the isoform difference in the resolved active site. An unshifted one would say the active
sites do not differ in this respect, and would point instead at the region where the structures
plainly do differ: GAD65's catalytic loop is disordered --- nine residues of it cannot be placed at
all --- where GAD67's is fully built, and GAD65's mean atomic displacement parameter is twice
GAD67's at the same resolution. A clamped calculation cannot represent disorder; but by excluding the
alternative it can say that disorder is where the difference must live.

\begin{thebibliography}{9}
\bibitem{fenalti} G.~Fenalti et al., \emph{GABA production by glutamic acid decarboxylase is
regulated by a dynamic catalytic loop}, Nat.\ Struct.\ Mol.\ Biol.\ \textbf{14} (2007) 280--286.
PDB entries 2OKK (GAD65) and 2OKJ (GAD67).
\bibitem{realqm} C.~Johnson, \emph{Real Quantum Mechanics}, and the RealQM solver,
\url{https://github.com/Claes542/RealMolecule}. Runs in a WebGPU browser with no installation.
\bibitem{dunathan} H.~C.~Dunathan, \emph{Conformation and reaction specificity in pyridoxal
phosphate enzymes}, Proc.\ Natl.\ Acad.\ Sci.\ USA \textbf{55} (1966) 712--716.
\end{thebibliography}

\end{document}
